\section{Proofs for Section~\ref{sec:split}}
\label{app:split}

As in Section~\ref{app:regions}, $B_\Hh$ is the unit ball of $\Hh$ and $\mathcal G_\alpha(P)$ the set of
trimming densities; the betting notation is that of Section~\ref{sec:betting}. We begin with two facts about separation in a fixed direction, which the other sections also use.

\subsection{Separation in a fixed direction}

The distance between two convex sets is the largest gap between them along a unit direction.

\begin{lemma}[Directional separation]
\label{lem:separation}
Let $A,B\subset\Hh$ be nonempty. For every $f\in B_\Hh$,
\begin{equation}
\dist(A,B)\ \ge\ \inf_{a\in A}\langle f,a\rangle_\Hh-\sup_{b\in B}\langle f,b\rangle_\Hh\ =:\ \mathrm{sep}(f;A,B).
\label{eq:sep}
\end{equation}
If $A,B$ are also convex and weakly compact and $\dist(A,B)>0$, the bound is attained: taking
$f^\star=(a^\star-b^\star)/\|a^\star-b^\star\|_\Hh$ for nearest points $a^\star\in A$, $b^\star\in B$ gives
$\mathrm{sep}(f^\star;A,B)=\dist(A,B)$.
\end{lemma}

\begin{proof}
For $a\in A$ and $b\in B$, the Cauchy--Schwarz inequality gives
$\mathrm{sep}(f;A,B)\le\langle f,a-b\rangle_\Hh\le\|a-b\|_\Hh$, and the infimum over $a,b$ yields \eqref{eq:sep}. For
the equality, nearest points exist, as the norm is weakly lower semicontinuous on the weakly compact $A\times B$, and
we check that $f^\star$ separates the sets exactly at them. First-order optimality of $a^\star$ for
$\min_{a\in A}\|a-b^\star\|_\Hh^2$ gives $\langle a-a^\star,a^\star-b^\star\rangle_\Hh\ge0$ for all $a\in A$, that is,
$\inf_{a\in A}\langle f^\star,a\rangle_\Hh=\langle f^\star,a^\star\rangle_\Hh$; symmetrically
$\sup_{b\in B}\langle f^\star,b\rangle_\Hh=\langle f^\star,b^\star\rangle_\Hh$. Subtracting,
$\mathrm{sep}(f^\star;A,B)=\langle f^\star,a^\star-b^\star\rangle_\Hh=\|a^\star-b^\star\|_\Hh=\dist(A,B)$.
\end{proof}

Since $\mathrm{sep}(0;A,B)=0$, for convex weakly compact $A,B$ the lemma gives
\begin{equation}
\dist(A,B)=\max_{f\in B_\Hh}\mathrm{sep}(f;A,B),
\label{eq:distmax}
\end{equation}
the maximum being attained at $f^\star$ when $\dist(A,B)>0$ and at $f=0$ otherwise. By Lemma~\ref{lem:attainment} and
Proposition~\ref{prop:georep} this applies to the population and to the empirical regions, and by Lemma~\ref{lem:cvar},
$\mathrm{sep}(f;\mathcal C_\alpha(P),\mathcal C_\alpha(Q))=-h_{\mathcal C_\alpha(P)}(-f)-h_{\mathcal C_\alpha(Q)}(f)$. For
empirical regions both terms are trimmed means of the values of $f$, which gives the structure of the empirical
trimming (Proposition~\ref{prop:bangbang}, proved in Section~\ref{app:trimming}).

\begin{lemma}[Support functions are extreme trimmed means]
\label{lem:trimmedmean}
Let $v_1,\dots,v_n$ be reals with order statistics $v_{(1)}\le\dots\le v_{(n)}$, let $\alpha\in[0,1)$ and
$c_\alpha=[n(1-\alpha)]^{-1}=\beta/n$. Then
\begin{gather*}
\max_{\substack{0\le w_i\le c_\alpha\\ \sum_iw_i=1}}\ \sum_iw_iv_i
\ =\ \sum_{i>n-n_{1-\alpha}}c_\alpha v_{(i)}+\big(1-n_{1-\alpha} c_\alpha\big)\,v_{(n-n_{1-\alpha})},\\
n_{1-\alpha}=\lfloor1/c_\alpha\rfloor=\lfloor n(1-\alpha)\rfloor,
\end{gather*}
where the final term is $0$ when $n_{1-\alpha}=n$ (at $\alpha=0$); the minimum is the same expression applied to
the smallest values. Moreover, every maximizer satisfies $w_i=c_\alpha$ when $v_i>v_{(n-n_{1-\alpha})}$ and $w_i=0$ when
$v_i<v_{(n-n_{1-\alpha})}$, the values equal to $v_{(n-n_{1-\alpha})}$ sharing the remaining mass $1-c_\alpha\#\{i:v_i>v_{(n-n_{1-\alpha})}\}$ in any way. So the maximizer is unique for distinct values.
\end{lemma}

\begin{proof}
Moving weight from a smaller value to a larger one below its cap does not decrease the objective, so an optimum loads
the cap $c_\alpha$ onto the largest values until the mass runs out, which gives the value. For the structure of the
maximizers, let $n_{1-\alpha}<n$ (otherwise the feasible set is a point) and suppose that a maximizer $w$ has
$w_i<c_\alpha$ for some $v_i>v_{(n-n_{1-\alpha})}$. At most $n_{1-\alpha}$ values are at least $v_i$, and they carry less than
$n_{1-\alpha} c_\alpha\le1$, so some $j$ with $v_j<v_i$ has $w_j>0$; moving weight from $j$ to $i$ increases the objective,
a contradiction. The case $v_i<v_{(n-n_{1-\alpha})}$ with $w_i>0$ is symmetric. The objective depends on the mass at the
value $v_{(n-n_{1-\alpha})}$ only through its total, and the minimum follows by applying all this to $-v_1,\dots,-v_n$.
\end{proof}

\subsection{Proof of Lemma~\ref{lem:ramp}}

\begin{proof}
Part (i) is the variational form of trimmed means of Lemma~\ref{lem:cvar}, part (ii) an exact computation, and
part (iii) the layer-cake formula. For the inequality in (i), let $P'\in R_\alpha(P)$ and $g=dP'/dP\in[0,\beta]$, so
that $\E_Pg=1$. Then $\E_{P'}f=t+\E_P[g(f-t)]\ge t-\E_P[g(t-f)_+]\ge t-\beta\E_P(t-f)_+$, and symmetrically $\E_{Q'}f\le u+\beta\E_Q(f-u)_+$ for $Q'\in R_\alpha(Q)$; subtracting and optimizing over $P'$ and $Q'$
gives the inequality. For the equality, Lemma~\ref{lem:cvar} applied to the bounded function $-f$, with its variable
written $-t$, gives $\inf_{P'\in R_\alpha(P)}\E_{P'}f=-\sup_{g\in\mathcal G_\alpha(P)}\E_P[-gf]=\sup_{t\in\R}\{t-\beta\E_P(t-f)_+\}$,
attained at $t=-t'$ for an $\alpha$-quantile $t'$ of $-f(X)$, that is, at a $(1-\alpha)$-quantile of $f(X)$. Applied to
$f$ under $Q$, the lemma gives $\sup_{Q'\in R_\alpha(Q)}\E_{Q'}f=\inf_{u\in\R}\{u+\beta\E_Q(f-u)_+\}$, attained at an
$\alpha$-quantile of $f(Y)$. For $f\in B_\Hh$, the reproducing property and Proposition~\ref{prop:georep} turn the two
extreme trimmed means into $\inf_{a\in\mathcal C_\alpha(P)}\langle f,a\rangle_\Hh$ and
$\sup_{b\in\mathcal C_\alpha(Q)}\langle f,b\rangle_\Hh$, whose difference is the separation \eqref{eq:sep}.

For (ii), clipping makes both hinges linear in the ramp: since $\tilde f\le t$ and $\tilde f\ge u$,
\[
t-\beta(t-\tilde f)_+=t-\beta(t-u)(1-\rp),
\qquad
u+\beta(\tilde f-u)_+=u+\beta(t-u)\rp ,
\]
and taking expectations and using $1-\beta=-\alpha\beta$ gives the identity. For the inequality, pointwise
$(t-\tilde f)_+=(t-f)_+-(u-f)_+$ and $(\tilde f-u)_+=(f-u)_+-(f-t)_+$, so
$\theta(\tilde f;t,u)-\theta(f;t,u)=\beta\E_P(u-f)_++\beta\E_Q(f-t)_+\ge0$.

For (iii), if $\rp$ takes values in $[0,1]$, the layer-cake formula and the definition of total variation give
$\E_P\rp-\E_Q\rp=\int_0^1\{P(\rp>r)-Q(\rp>r)\}\,dr\le\TV(P,Q)$.
\end{proof}

\subsection{The betting test in full}
\label{sec:bettingdetails}

We describe the test of Section~\ref{sec:betting} for any batch fractions $0=\pi_0<\pi_1<\dots<\pi_{B+1}=1$ fixed in advance;
Section~\ref{sec:betting} uses $\pi=(0,\tfrac14,\tfrac12,1)$, and the single split $(0,\tfrac12,1)$ is the case $B=1$. Each
sample is put in a random order drawn independently of the data and cut into batches $0,\dots,B$ at these fractions of
its size, rounded to the nearest integers. Batch $0$ is used
only for fitting, and batch $b\ge1$ is bet on with a function $\rp_b$ with values in $[0,1]$ fitted on batches
$0,\dots,b-1$, or not at all.

Before betting on batch $b$ we solve \eqref{eq:qp} on batches $0,\dots,b-1$ and form the direction
$\widehat f_b$ of the witness \eqref{eq:witness}. Its cut points are $\widehat t_b$, a genuine empirical
$(1-\alpha)$-quantile of its values at the $X$'s of those batches, and $\widehat u_b$, a genuine empirical
$\alpha$-quantile of its values at their $Y$'s, each value computed with the point's own term removed from
\eqref{eq:witness}. If $\widehat
t_b>\widehat u_b$, $\rp_b$ is the ramp of $\widehat f_b$ between $\widehat u_b$ and $\widehat t_b$ of
Lemma~\ref{lem:ramp}(ii); otherwise there is no bet on batch $b$.

Pairs are formed within each batch. Let batch $b$ hold $n_{[b]}$ observations $X$ and $m_{[b]}$ observations $Y$. If $n_{[b]}\le m_{[b]}$,
its $Y$'s are cut, in their fixed order, into $n_{[b]}$ consecutive blocks whose sizes differ by at most one, and the
$X$'s are paired one to one with the blocks, an $X$ and its block $\mathcal J$ giving
$Z=\rp_b(X)-|\mathcal J|^{-1}\sum_{j\in\mathcal J}\rp_b(Y_j)$. If $n_{[b]}>m_{[b]}$, the roles are exchanged: each $Y$ is
paired with a block $\mathcal J$ of $X$'s, giving $Z=|\mathcal J|^{-1}\sum_{i\in\mathcal J}\rp_b(X_i)-\rp_b(Y)$. Batch $b$
thus gives $N_b=\min(n_{[b]},m_{[b]})$ pairs, and the pairs of all batches are numbered $k=1,\dots,N$ in the order of the
batches, with $N=\sum_bN_b$ fixed by the sample sizes. Each $Z_k$ takes values in $[-1,1]$.

The wealth after $k$ pairs is $K_k=\prod_{l\le k}\{1+\lambda_l(Z_l-\alpha)\}$, with
$K_0=1$, and both the function bet on and the fractions must be predictable. A fitting rule is predictable if each $\rp_b$ is a jointly measurable function of batches $0,\dots,b-1$ and of the
point, with values in $[0,1]$; the ramp above is one, and so are the predicted labels of a classifier trained on batches
$0,\dots,b-1$. Fractions are predictable if, for pair $k$ in batch $b$, $\lambda_k\in[0,\lambda_{\max}/(1+\alpha)]$ for a
fixed $\lambda_{\max}\in(0,1]$, $\lambda_k$ is a jointly measurable function of batches $0,\dots,b-1$ and of the earlier
$Z$'s, and $\lambda_k=0$ on a batch without a ramp. Theorem~\ref{thm:betting} holds, with the same statement, for every
choice of batch fractions $\pi$; we prove it in this generality.

\begin{proof}[Proof of Theorem~\ref{thm:betting}]
Under the null every bet has a nonpositive expected gain, so the wealth is a nonnegative supermartingale, and such a
process rarely reaches $1/\delta$. To set up the filtration, for pair $k$ in batch $b$ let $\mathcal F_{k-1}$ be generated
by batches $0,\dots,b-1$ and the earlier pairs of batch $b$. These $\sigma$-fields increase with $k$, $\rp_b$ and
$\lambda_k$ are $\mathcal F_{k-1}$-measurable, and the observations of pair $k$ are independent of $\mathcal F_{k-1}$.
So, given $\mathcal F_{k-1}$, $Z_k$ takes values in $[-1,1]$ and has mean
$\Delta_k=\E_P\rp_b-\E_Q\rp_b\le\TV(P,Q)\le\alpha$, by Lemma~\ref{lem:ramp}(iii) and the null. Consequently
$\E[K_k\mid\mathcal F_{k-1}]=K_{k-1}\{1+\lambda_k(\Delta_k-\alpha)\}\le K_{k-1}$, while
$\lambda_k(Z_k-\alpha)\ge-\lambda_k(1+\alpha)\ge-\lambda_{\max}$ gives $K_k\ge(1-\lambda_{\max})K_{k-1}\ge0$. So $(K_k)$ is
a nonnegative supermartingale with $K_0=1$.

Optional stopping at $T=\min\{k\le N:K_k\ge1/\delta\}\wedge N$, a bounded time as $N$ is fixed in advance, gives
$\E[K_T\mid\text{batch }0]\le1$. As $\{\max_{k\le N}K_k\ge1/\delta\}\subseteq\{K_T\ge1/\delta\}$, Markov's inequality
gives the bound, which is Ville's maximal inequality \citep{ville1939}. Integrating over batch $0$ ends the proof.
\end{proof}

The fractions we use are
\begin{equation}
\lambda_k=\min\Big\{\frac{\lambda_{\max}}{1+\alpha},\ \sqrt{\frac{2\log(1/\delta)}{N\widehat v^2_{k-1}}}\Big\},
\qquad \lambda_{\max}=\tfrac12,
\label{eq:fractions}
\end{equation}
adapted from \citet[Remark~1]{waudbysmith2024estimating}, where
$\widehat v^2_{k-1}=(w_0v_0^2+\sum_{i<k}(Z_i-\alpha-\widehat\mu_{i-1})^2)/(w_0+k-1)$,
$\widehat\mu_{i-1}=\sum_{l<i}(Z_l-\alpha)/(w_0+i-1)$, $w_0=20$, $v_0^2\le4$ is an estimate of the variance of one
$Z_k$ from the batches before that of $Z_k$ (Algorithm~\ref{alg:split} gives ours), the sums and the counts $k-1$ and
$i-1$ run over the pairs actually bet on before $Z_k$, in its batch and in the earlier ones, so that they carry over
from batch $1$ to batch $2$ while $v_0^2$ is that of the batch of $Z_k$, and $\widehat v^2_{k-1}$ is floored at a constant $v^2_{\min}>0$. Our runs use
$v^2_{\min}=10^{-8}$, a guard against division by zero, for which the constants of Proposition~\ref{prop:loconsistency}
and Theorems~\ref{thm:lorate} and~\ref{thm:laplacerate}, of order $1/v^2_{\min}$, are not usable and those results
describe the trend (Section~\ref{sec:lower} of the main text). Any other predictable choice in
$[0,\lambda_{\max}/(1+\alpha)]$ is equally valid, and this one comes close to the best constant fraction for large $N$.

\begin{algorithm}[!tb]
\caption{The split-sample lower bound $\widehat\alpha$, with the defaults of our runs}
\label{alg:split}
\begin{algorithmic}[1]
\Require samples $X_{1:n}$, $Y_{1:m}$; level $\delta$; a grid $\alpha_1<\dots<\alpha_J$, fixed in advance or chosen from
batch $0$ (Remark~\ref{rem:batch0grid})
\State Put each sample in a random order drawn independently of the data and cut it into batch $0$ (first quarter),
batch $1$ (second quarter) and batch $2$ (second half)
\State Pair within each batch $b$: if $n_{[b]}\le m_{[b]}$, cut its $Y$'s in order into $n_{[b]}$ consecutive blocks of
sizes differing by at most one and pair the $i$-th $X$ with the $i$-th block, and the reverse otherwise; $N$ is the
total number of pairs
\For{$j=1,\dots,J$, with $\alpha=\alpha_j$}
\For{$b=1,2$, fitting on batches $0,\dots,b-1$, of sizes $n_{<b}$, $m_{<b}$}
\State Kernel $e^{-\|x-y\|/\sigma}$, with $\sigma$ the median pairwise distance of the fitting points (simulations) or
of batch $0$ (real data, both batches; RHC standardized by the moments of batch $0$), on at most $2{,}000$ points
\State Minimizer $(w,v)$ of \eqref{eq:qp} at $\alpha$ by accelerated projected gradient with restarts and an exact
projection onto the capped simplex, warm-started from the previous level, stopped when the objective changes by at most
a relative $10^{-9}$ over ten iterations or after $1000$ iterations ($10^{-10}$ and $2000$ on real data)
\State Witness values at the fitting points without their own term, $\tilde h_i=\widehat f(X_i)-w_ik(X_i,X_i)/\widehat D_\alpha$
and $\tilde h'_j=\widehat f(Y_j)+v_jk(Y_j,Y_j)/\widehat D_\alpha$
\State Cut points $\widehat t_b=\tilde h_{(\lceil n_{<b}(1-\alpha)\rceil)}$ and $\widehat u_b=\tilde h'_{(\lfloor m_{<b}\alpha\rfloor+1)}$,
order statistics
\If{$\widehat t_b\le\widehat u_b$} skip batch $b$: $\lambda_k=0$ on its pairs
\Else{} $\rp_b$ is the ramp of $\widehat f_b$ between $\widehat u_b$ and $\widehat t_b$, each pair gives
$Z=\rp_b(X)-\overline{\rp_b(Y)}$ (or the mirror), and $v_0^2=\max\{s_X^2+s_Y^2n_{[b]}/m_{[b]},10^{-4}\}$ (mirrored
if $n_{[b]}>m_{[b]}$), with $s_X^2,s_Y^2$ the variances of $\rp_b$ at the fitting $X$'s and $Y$'s
\EndIf
\EndFor
\State Bet: $K_0=1$, $K_k=K_{k-1}\{1+\lambda_k(Z_k-\alpha)\}$, with \eqref{eq:fractions}, $\lambda_{\max}=\tfrac12$,
$w_0=20$ and $v^2_{\min}=10^{-8}$
\If{$\max_{k\le N}K_k<1/\delta$} stop: $H_{0,\alpha_j}$ is not rejected \EndIf
\EndFor
\State \Return $\widehat\alpha$ of \eqref{eq:downclosed}; whether the reading stopped at a failed bet or with both batches
skipped; the trimming at $\widehat\alpha$, the minimizer of \eqref{eq:qp} on all the data closest to uniform weights
(Section~\ref{sec:computationdetails})
\end{algorithmic}
\end{algorithm}

Both bandwidth rules are predictable, each witness using a bandwidth computed from the batches before the one bet on,
but the consistency and rate results, which fix the kernel in advance, cover neither of them, nor random features
scaled by such a bandwidth. Section~\ref{app:experiments} gives the grids of our experiments; in practice we
recommend a grid of mesh $0.05$ fixed in advance, or a short grid around a guess at $\alpha^\star$ from batch $0$
(Remark~\ref{rem:batch0grid}).

The constant $\cstBetRatioEighty$ of Section~\ref{sec:betting} comes from the best constant fraction: $\log
K_N=\lambda S_1-\tfrac12\lambda^2S_2+O(\lambda^3N)$, with $S_1=\sum_k(Z_k-\alpha)$ and $S_2=\sum_k(Z_k-\alpha)^2$, is at
most about $S_1^2/(2S_2)$, so the test rejects roughly when $S_1/S_2^{1/2}\ge\sqrt{2\log(1/\delta)}$, an approximately
$N(h,1)$ ratio if $\E Z_k-\alpha=hsN^{-1/2}$, with $s^2=\operatorname{Var}Z_k$. Power $p$ then needs $h\approx\sqrt{2\log(1/\delta)}+z_p$, against
$z_{1-\delta}+z_p$ by the central limit theorem, and the squared ratio at $\delta=0.05$, $p=0.8$ is the ratio of sizes.

\subsection{Proof of Corollary~\ref{cor:lower}}

\begin{proof}
A bound that is too large requires the rejection of one particular true null, the first grid level at or
above $\TV(P,Q)$, so no union over the levels is needed. Let $E$ be the event that $H_{0,\widehat\alpha}$ is rejected and $\TV(P,Q)\le\widehat\alpha$; it
contains $\{\widehat\alpha>\TV(P,Q)\}$, since then $\widehat\alpha>0$ is a rejected grid level. On
$E$, $\widehat\alpha=\alpha_j$ with $\alpha_j\ge\TV(P,Q)$ and $H_{0,\alpha_i}$ rejected for every $i\le j$. The
smallest index $i_0$ with $\alpha_{i_0}\ge\TV(P,Q)$ is at most $j$ and does not depend on the data, because the grid is
fixed in advance, so $E$ lies in the event that the true null $H_{0,\alpha_{i_0}}$ is rejected, of probability at most
$\delta$ by Theorem~\ref{thm:betting}. No multiplicity correction is needed, and if no grid level is at least
$\TV(P,Q)$, $E$ is empty.
\end{proof}

\begin{remark}[A grid chosen from batch $0$]
\label{rem:batch0grid}
Corollary~\ref{cor:lower} holds, with the same proof, when the grid $\alpha_1<\dots<\alpha_J$ is a measurable function
of batch $0$ instead of being fixed in advance. Given batch $0$ the grid is fixed, so the smallest index $i_0$ with
$\alpha_{i_0}\ge\TV(P,Q)$ is a function of batch $0$, and the conditional bound of Theorem~\ref{thm:betting} gives
$\Pp(H_{0,\alpha_{i_0}}\text{ rejected}\mid\text{batch }0)\le\delta$; integrating over batch $0$ gives the claim. For
instance, a bisection on batch $0$ alone, with any statistic, may place a short grid around a guess at $\alpha^\star$,
which the betting batches then read upward.
\end{remark}

\subsection{Proof of Proposition~\ref{prop:loconsistency}}

We allow any batch fractions fixed in advance. For $b\ge1$ let $n_{<b},m_{<b}$ be the sizes of batches $0,\dots,b-1$,
$r_b=1/n_{<b}+1/m_{<b}\le r_1$, and $\Gamma_b,\Gamma_b^-$ the suprema $\Gamma,\Gamma^-$ of the proof of
Theorem~\ref{thm:twosided} on these batches, free of $\alpha$. At a level $\alpha$ the witness fitted on them takes
values $h_i,h'_j$ there, and $\tilde h_i=h_i-w_ik(X_i,X_i)/\widehat D$ and
$\tilde h'_j=h'_j+v_jk(Y_j,Y_j)/\widehat D$ with the own term removed. With $\iota^-(n,p)=\max(1,\lceil np\rceil-1)$ and
$\iota^+(n,p)=\min(n,\lfloor np\rfloor+2)$, the cut points satisfy
\begin{equation}
\tilde h_{(\iota^-(n_{<b},1-\alpha))}\le\widehat t_b\le\tilde h_{(\iota^+(n_{<b},1-\alpha))},\qquad
\tilde h'_{(\iota^-(m_{<b},\alpha))}\le\widehat u_b\le\tilde h'_{(\iota^+(m_{<b},\alpha))},
\label{eq:cutrank}
\end{equation}
as do the empirical quantiles in the range of the values and their linear interpolations.

\begin{lemma}[The population ramp]
\label{lem:popramp}
Under (A1) and (A2), if $D_\alpha>0$, the ramp $\rp^\star$ of the normalized difference $f^\star$ of nearest points of
$\mathcal C_\alpha(P)$ and $\mathcal C_\alpha(Q)$, between the largest $\alpha$-quantile $u^\star$ of $f^\star(Y)$ and the
smallest $(1-\alpha)$-quantile $t^\star$ of $f^\star(X)$, has $D_\alpha\le t^\star-u^\star\le \Delta_\phi$ and
$\E_P\rp^\star-\E_Q\rp^\star-\alpha\ge D_\alpha/(\beta \Delta_\phi)$.
\end{lemma}

\begin{proof}
The right derivatives $1-\beta P(f^\star\le t)$ and $1-\beta Q(f^\star>u)$ make $t^\star$ maximize
$t-\beta\E_P(t-f^\star)_+$ and $u^\star$ minimize $u+\beta\E_Q(f^\star-u)_+$. So by Lemmas~\ref{lem:ramp}(i)
and~\ref{lem:separation} the cut points attain the separation, and since $\theta(f;t,u)\le t-u$,
$\Delta_\phi\ge\sup f^\star-\inf f^\star\ge t^\star-u^\star\ge\theta(f^\star;t^\star,u^\star)=D_\alpha$, the first
inequality because $f^\star$ has unit norm. Lemma~\ref{lem:ramp}(ii) then gives
$\beta(t^\star-u^\star)(\E_P\rp^\star-\E_Q\rp^\star-\alpha)\ge D_\alpha$, and $t^\star-u^\star\le\Delta_\phi$ gives the
excess.
\end{proof}

\begin{proof}[Proof of Proposition~\ref{prop:loconsistency}]
Below $\alpha^\star$ the population ramp has a fixed positive excess over $\alpha$ (Lemma~\ref{lem:popramp}). We show
that the fitted ramp keeps part of this excess on an event that does not depend on the level, that betting turns a fixed excess
into a rejection, and then apply both to the grid.

We begin with the fitted ramp at a fixed level $\alpha$ and batch $b$, writing $n,m$ for $n_{<b},m_{<b}$, and assume
$\beta\Gamma_b^-<D_\alpha$. By the deterministic reduction in the proof of Theorem~\ref{thm:twosided},
$\widehat D\ge D_\alpha-\beta\Gamma_b^->0$. Let $\theta_n$ be the criterion $\theta$ of Section~\ref{sec:ramp} with
$P_n,Q_m$ in place of $P,Q$. In the fitted direction $\widehat f$ it separates into a concave function of the first cut
and a convex function of the second,
\[
\theta_n(\widehat f;t,u)=T(t)-U(u),\quad T(t)=t-\frac\beta n\sum_i(t-h_i)_+,\quad U(u)=u+\frac\beta m\sum_j(h'_j-u)_+,
\]
and $\max T-\min U=\widehat D$ by Lemmas~\ref{lem:separation} and~\ref{lem:ramp}(i) applied to $P_n,Q_m$. We first show
that the cut points lose little against the best cuts. The concave $T$ has right derivative $1-\beta F_n$, with $F_n(x)=\#\{i:h_i\le x\}/n$, so it
is $\beta$-Lipschitz and maximal at $h_{(\lceil n(1-\alpha)\rceil)}$. Between $h_{(\iota^-)}$ and $h_{(\iota^+)}$ we have
$\iota^-/n\le F_n\le(\iota^+-1)/n$, within $1/n$ of $1-\alpha$, so the slope of $T$ there is at most $\beta/n$; as the
values span at most $\Delta_\phi$, $T\ge\max T-\beta\Delta_\phi/n$ on that interval. The own terms shift the values only
a little: the caps $\beta/n$ and $k(x,x)\le\kappa^2$ give $0\le h_i-\tilde h_i\le\beta\kappa^2/(n\widehat D)$, the order
statistics move by at most as much, and concavity, with the $\beta$-Lipschitz bound, gives
\[
T(\widehat t_b)\ge\max T-\frac{\beta\Delta_\phi}n-\frac{\beta^2\kappa^2}{n\widehat D} .
\]
The convex $U$, minimal at $h'_{(\lfloor m\alpha\rfloor+1)}$, is treated alike, with $m$ in place of $n$. Hence
$\theta_n(\widehat f;\widehat t_b,\widehat u_b)\ge\widehat D-\beta\Delta_\phi r_b-\beta^2\kappa^2r_b/\widehat D$.

Next we pass from the empirical criterion to the population one. For all $(f,t,u)$,
$\theta-\theta_n=-\beta[(P-P_n)(t-f)_++(Q-Q_m)(f-u)_+]\ge-\beta\Gamma_b$, the centered terms being constant in $t$
and $u$ outside $[-\kappa,\kappa]$. As $x\mapsto x-c/x$ increases and $\widehat D\ge D_\alpha-\beta\Gamma_b^-$, we may
replace $\widehat D$ by $D_\alpha-\beta\Gamma_b^-$. Finally $\theta(f;t,u)\le t-u$, and $\widehat t_b\le\max_ih_i$,
$\widehat u_b\ge\min_jh'_j$, values of a unit vector, are at most $\Delta_\phi$ apart. Together,
\begin{equation}
\Delta_\phi\ge\widehat t_b-\widehat u_b\ge\theta(\widehat f;\widehat t_b,\widehat u_b)\ge\Theta_b=D_\alpha-\beta(\Gamma_b
+\Gamma_b^-)-\beta \Delta_\phi r_b-\frac{\beta^2\kappa^2r_b}{D_\alpha-\beta\Gamma_b^-},
\label{eq:fittedramp}
\end{equation}
and if $\Theta_b>0$, Lemma~\ref{lem:ramp}(ii) under $P,Q$, given batches $0,\dots,b-1$, yields a ramp $\rp_b$ with
$\E_P\rp_b-\E_Q\rp_b-\alpha\ge\Theta_b/(\beta \Delta_\phi)$. The terms of $\Theta_b$ after $D_\alpha$ are the costs of the
fluctuation, of the rank tolerance and of the own terms.

The second step concerns the bets: a fixed excess is turned into a rejection, and a strict null is rarely rejected.
Recall that $\Delta_k$ is the conditional mean of $Z_k$ in the proof of Theorem~\ref{thm:betting}, and put
$L_\delta=\log(1/\delta)$ and
\[
\varepsilon_N(\varrho)=\exp\Big\{-\frac{v^2_{\min}}{4L_\delta}\Big(\varrho\sqrt{NL_\delta/2}-L_\delta
-\frac{16L_\delta}{v^2_{\min}}\Big)_+^2\Big\}.
\]
We show that for $N\ge2L_\delta/\lambda_{\max}^2$,
\begin{equation}
\Pp\big(\mathcal E\cap\{\max_{k\le N}K_k<1/\delta\}\big)\le\varepsilon_N(\varrho),\qquad
\Pp\big(\max_{k\le N}K_k\ge1/\delta\big)\le e^{-\zeta v_{\min}(2NL_\delta)^{1/2}},
\label{eq:betbounds}
\end{equation}
the first for any event $\mathcal E$ on which every batch has a ramp and $\Delta_k\ge\alpha+\varrho>\alpha$ for all $k$,
the second if $\TV(P,Q)\le\alpha-\zeta<\alpha$.

Two bounds on the fractions are used. Each square in \eqref{eq:fractions} is at most $4$, as $Z_i-\alpha$ and
$\widehat\mu_{i-1}$ lie in $[-1-\alpha,1-\alpha]$, so for $v^2_{\min}\le4$ the running variance lies in
$[v^2_{\min},4]$ and
\[
\Big(\frac{L_\delta}{2N}\Big)^{1/2}\ \le\ \lambda_k\ \le\ \overline\lambda_N
=\min\Big\{\frac{\lambda_{\max}}{1+\alpha},\Big(\frac{2L_\delta}{Nv^2_{\min}}\Big)^{1/2}\Big\},
\]
the lower bound holding on batches with a ramp, where it uses $N\ge2L_\delta/\lambda_{\max}^2$. The centered gains
$M_k=\sum_{l\le k}\lambda_l(Z_l-\Delta_l)$ have increments in intervals of length $2\lambda_l$, so by Hoeffding's
lemma \citep{hoeffding1963}
\[
W_k(s)=\exp\Big(sM_k-\frac{s^2}2\sum_{l\le k}\lambda_l^2\Big)
\]
is a nonnegative supermartingale for every real $s$.

On $\mathcal E$ the wealth grows. The inequality $\log(1+y)\ge y-2y^2$ for $y\ge-\frac12$, the lower bound on
$\lambda_k$ and $\Delta_k-\alpha\ge\varrho$ give
\[
\log K_N\ \ge\ \varrho\Big(\frac{NL_\delta}2\Big)^{1/2}+M_N-\frac{4L_\delta(1+\alpha)^2}{v^2_{\min}},
\]
the last term from $2\sum_k\lambda_k^2(Z_k-\alpha)^2\le2N\overline\lambda_N^{\,2}(1+\alpha)^2$. As $(1+\alpha)^2\le4$,
failing to reject on $\mathcal E$ forces $M_N\le-x_N$, with
$x_N=(\varrho\sqrt{NL_\delta/2}-L_\delta-16L_\delta/v^2_{\min})_+$. Put $V=2L_\delta/v^2_{\min}\ge\sum_l\lambda_l^2$.
Markov's inequality for $W_N(-x_N/V)$, which is at least $e^{x_N^2/(2V)}$ when $M_N\le-x_N$, bounds the probability
by $e^{-x_N^2/(2V)}=\varepsilon_N(\varrho)$, without conditioning on $\mathcal E$.

Under the strict null the wealth shrinks. Now $\log(1+y)\le y$ and $\Delta_l\le\alpha-\zeta$ give
$\log K_k\le M_k-\zeta\sum_{l\le k}\lambda_l$. With $s=2\zeta/\overline\lambda_N$ and $\lambda_l\le\overline\lambda_N$,
the quadratic term of $W_k(s)$ is at most $\frac{s^2}2\overline\lambda_N\sum_{l\le k}\lambda_l=s\zeta\sum_{l\le k}
\lambda_l$, so $W_k(s)\ge K_k^{\,s}\ge\delta^{-s}$ wherever $K_k\ge1/\delta$. The maximal inequality in the proof of
Theorem~\ref{thm:betting} bounds the probability of this by $\delta^{2\zeta/\overline\lambda_N}$, which is at most the
second bound of \eqref{eq:betbounds} by the bound on $\overline\lambda_N$.

We now apply the two steps to the grid, starting with the upper side. If $\widehat\alpha>\alpha_{j^\star}$, then by
\eqref{eq:downclosed} the true null at the first grid level $\alpha_{i_0}\ge\alpha^\star$ was rejected, with probability
at most $\delta$ by Theorem~\ref{thm:betting}, and tending to zero if $\alpha^\star$ is off the grid, by
\eqref{eq:betbounds} with $\zeta=\alpha_{i_0}-\alpha^\star$.

On the lower side every grid level below $\alpha^\star$ must be rejected. Let $\alpha_{j^\star}$ be, as in the
statement, the largest grid level below $\alpha^\star$, $\underline D=D_{\alpha_{j^\star}}>0$ by \eqref{eq:zeroeq}, $\bar\beta=(1-\alpha_{j^\star})^{-1}$ and
$\beta_j=(1-\alpha_j)^{-1}$, so that $D_{\alpha_j}\ge \underline D$ and $\beta_j\le\bar\beta$ for $j\le j^\star$. Let $\mathcal E_{\mathrm{fit}}$ be the
event that $\Gamma_b,\Gamma_b^-\le \underline D/(8\bar\beta)$ for all $b$, and let the samples be large enough
that $\bar\beta\Delta_\phi r_1\le \underline D/8$ and $\bar\beta^2\kappa^2r_1\le7\underline D^2/64$. On
$\mathcal E_{\mathrm{fit}}$, at every level $\alpha_j$ with $j\le j^\star$, the four terms of $\Theta_b$ are at least
$\underline D$, $-\underline D/4$, $-\underline D/8$ and $-\underline D/8$, the last because the denominator is at
least $7\underline D/8$; so every batch has a ramp with excess at least $\varrho_*=\underline D/(2\bar\beta\Delta_\phi)$.
By Lemma~\ref{lem:hinge}, as in the proof of Theorem~\ref{thm:twosided},
$\E\Gamma_b,\E\Gamma_b^-\le\mu_1=2(1+\sqrt2)\kappa\sqrt{r_1}$, and a deviation $y$ above the mean has probability
at most $e^{-2y^2/(\Delta_\phi^2r_1)}$. So by \eqref{eq:downclosed} and \eqref{eq:betbounds},
\[
\Pp(\widehat\alpha<\alpha_{j^\star})\le\Pp(\mathcal E_{\mathrm{fit}}^c)+J\varepsilon_N(\varrho_*)\le
2Be^{-2(\underline D/(8\bar\beta)-\mu_1)_+^2/(\Delta_\phi^2r_1)}+J\varepsilon_N(\varrho_*)\to0 .
\]

Finer grids need one more step. For $\epsilon\in(0,\alpha^\star)$, the lower side with
$\underline D=D_{\alpha^\star-\epsilon/2}$ and $\bar\beta=(1-\alpha^\star+\epsilon/2)^{-1}$ rejects all grid levels in
$[0,\alpha^\star-\epsilon/2]$ with probability tending to one, as $\log J=o(N)$. Then
$\widehat\alpha\ge\alpha^\star-\epsilon/2-2h_{\mathcal G}\ge\alpha^\star-\epsilon$ once the mesh
$h_{\mathcal G}=\sup_{a\in[0,1)}\dist(a,\mathcal G\cup\{0\})$ of the grid $\mathcal G$, a two-sided mesh over $[0,1)$,
not the $\Delta_{\mathcal G}$ of Theorem~\ref{thm:lorate}, is at most $\epsilon/4$, while
$\widehat\alpha>\alpha^\star+\epsilon$ needs the rejection of the first grid level above $\alpha^\star+\epsilon/2$, of
probability at most $e^{-\epsilon v_{\min}(NL_\delta/2)^{1/2}}$ by \eqref{eq:betbounds}.
\end{proof}

The constants grow like $1/v^2_{\min}$: with a floor that shrinks with $N$, $\overline\lambda_N$ tends to
$\lambda_{\max}/(1+\alpha)$, and the second bound of \eqref{eq:betbounds} only to
$\delta^{2\zeta(1+\alpha)/\lambda_{\max}}$. The proof also allows an inexact minimizer of \eqref{eq:qp}: if its
objective is $d^2$ and its Frank--Wolfe gap is at most $\rho d^2$ with $\rho<1$, the gap equals $2d^2-2ds$ for the
empirical separation $s$ of its witness, so $s\ge(1-\rho/2)d$, and the argument goes through with
$(1-\rho/2)\underline D$ in place of $\underline D$.

\subsection{A witness fitted at $\alpha=0$}
\label{sec:alphazero}

Consistency needs the witness fitted at the level tested. A single witness fitted at $\alpha=0$, with direction
$f_0=(\mu_P-\mu_Q)/\MMD(P,Q)$ as in a test of $P=Q$, can keep $\widehat\alpha$ away from $\alpha^\star$. We show that
its ramps certify at most a one-sided Kolmogorov distance of the values of $f_0$, and then give laws for which this
distance is smaller than $\TV(P,Q)$. A ramp $\rp$ of $f\in B_\Hh$ with cuts $t>u$ has
$\{\rp>r\}=\{f>u+r(t-u)\}$ for $r\in[0,1)$, so by the layer-cake formula
\[
\E_P\rp-\E_Q\rp\ \le\ \mathrm{KS}^+(f)=\sup_x\{P(f(X)>x)-Q(f(Y)>x)\}.
\]
Witnesses $\widehat f_{0,b}$ fitted at $\alpha=0$ have uniform weights and tend to $f_0$ uniformly, almost surely.
Multiplying the fractions by $\mathbf 1\{\mathrm{KS}^+(\widehat f_{0,b})\le\alpha\}$, which is predictable, gives a
test that has $\Delta_k\le\alpha$ wherever its fraction is positive. By the proof of Theorem~\ref{thm:betting}, this
test has level $\delta$ even when $H_{0,\alpha}$ is false, and it has the same wealth as the original test when
$\max_b\mathrm{KS}^+(\widehat f_{0,b})\le\alpha$. Hence the original test rejects $H_{0,\alpha}$ with probability at
most $\delta+\Pp(\max_b\mathrm{KS}^+(\widehat f_{0,b})>\alpha)$.

For the laws, take the Gaussian kernel with $\sigma=\loExSigma$ on $\R$, let
$P$ put mass $\loExPzero$ at $0$, $\loExCommon$ at $\loExCommonLoc$ and $\loExSpread$ on
$\loExSpreadLo,\loExSpreadNext,\dots,\loExSpreadHi$ in proportion to $1,2,\dots,\loExSpreadK$, and $Q$ put $\loExQhalf$
at $\loExQloc$, $\loExQfar$ at $\loExFar$ and $\loExCommon$ at $\loExCommonLoc$, so $\TV(P,Q)=\loExTV$. Up to
$O(e^{-\loExTail})$, $\mu_P-\mu_Q$ is $\loExPzero-\loExQhalf e^{-1/8}=\loExFzero$ at $0$, $\loExPzero
e^{-1/8}-\loExQhalf=\loExFhalf$ at $\loExQloc$, increasing and at most $\loExFspread$ on the spread points, $0$ at
$\loExCommonLoc$ and $-\loExQfar$ at $\loExFar$. So the $\loExAtoms$ values of $f_0$ are distinct, its up-sets have
$P-Q$ equal to $\loExPzero$, $\loExCumTwo$, at most $\loExCumSpread$, $\loExCumSpread$ and $0$, and
$\mathrm{KS}^+(\widehat f_{0,b})=\mathrm{KS}^+(f_0)=\loExKS$ once $\widehat f_{0,b}$ is uniformly within half the smallest
gap of these values, so every level $\alpha\ge\loExKS$ is rejected with probability at most $\delta+o(1)$, the false
nulls below $\TV(P,Q)$ included.
